\documentclass[11pt]{article}

\usepackage[margin=1in]{geometry}
\usepackage{amsmath,amssymb,amsthm}
\usepackage{mathtools}
\usepackage[hidelinks]{hyperref}

\newtheorem{theorem}{Theorem}
\newtheorem{lemma}{Lemma}
\newtheorem{proposition}{Proposition}
\newtheorem{corollary}{Corollary}
\theoremstyle{definition}
\newtheorem{definition}{Definition}
\newtheorem{assumption}{Assumption}
\theoremstyle{remark}
\newtheorem{remark}{Remark}

\DeclareMathOperator{\Poisson}{Poisson}
\DeclareMathOperator{\Binomial}{Binomial}
\DeclareMathOperator{\Var}{Var}
\DeclareMathOperator{\dTV}{d_{\mathrm{TV}}}
\newcommand{\Pgen}{P_{\mathrm{gen}}}
\newcommand{\Pzero}{P_{0}}
\newcommand{\ball}[2]{B_{#1}(#2)}
\newcommand{\R}{\mathbb{R}}
\newcommand{\E}{\mathbb{E}}
\newcommand{\Prob}{\mathbb{P}}
\newcommand{\one}{\mathbf{1}}

\title{A control-calibrated E-value for fuzzy TCR sequence search\\ over biologically redundant reference sets}
\author{seqtree --- technical appendix}
\date{}

\begin{document}
\maketitle

\begin{abstract}
We derive a BLAST-style E-value~\cite{Karlin1990,Altschul1990} for ``hits'' returned by fuzzy
search over T-cell receptor (TCR) CDR3 sequences, adapted to the defining difficulty of immune
repertoires: the reference set is highly \emph{redundant}, and the redundancy is biological
(convergent V(D)J recombination, public clones, clonal expansion) rather than statistical noise.
The classical Karlin--Altschul theory assumes a database of independent, identically distributed
letters; under that null, redundancy-driven near-matches are absurdly significant. We replace the
i.i.d.-letter null with an \emph{empirical background} $\Pzero$ estimated from a matched control
repertoire, retain the Poisson/Gumbel limit superstructure with an explicit non-asymptotic error
bound (Chen--Stein / Le~Cam~\cite{Arratia1990,LeCam1960}), and handle clonal over-dispersion by
collapsing to unique clonotypes. The resulting E-value is automatically deflated for hits that the
generation process alone explains and is large only for antigen-driven convergence. This puts the TCRNET approach---counting
sequence neighbours against a real-world control repertoire, first introduced by
Ritvo~et~al.~\cite{Ritvo2018} and formalized as an annotation framework
by~\cite{Pogorelyy2019}---on a rigorous, finite-sample footing, and we show the classical
Karlin--Altschul E-value is its product-measure, ungapped special case.
\end{abstract}

\section{Introduction: the redundancy problem}

Given a query CDR3 $q$ and a target set $D$ (e.g.\ VDJdb), fuzzy search returns the neighbours of
$q$ within a fixed scope/budget $\theta$. We want a significance value for such hits. BLAST answers
this for protein search with the Karlin--Altschul E-value
\begin{equation}\label{eq:ka}
  E = K\,m\,n\,e^{-\lambda^\ast S},
\end{equation}
where $m$ is the query length and $n$ the total database length (both in residues), $S$ is the
alignment score of the hit under a substitution matrix with entries $s_{ij}$, $\lambda^\ast$ is the
unique positive root of $\sum_{ij} p_i p_j e^{\lambda^\ast s_{ij}} = 1$ (the natural scale that turns
scores into log-probabilities for i.i.d.\ letters with background frequencies $p_i$), and $K>0$ is a
prefactor---the ``clumping'' or edge-effect constant---fixed by the score distribution. $E$ is the
expected number of distinct alignments scoring at least $S$ by chance; the number of such alignments
is asymptotically Poisson, so $\Pr(\text{at least one}) = 1-e^{-E}$~\cite{Karlin1990,Karlin1993}.
The whole construction rests on the database being a string of i.i.d.\ letters.

Immune repertoires violate the i.i.d.\ assumption catastrophically. CDR3s are produced by V(D)J
recombination, whose generation probability $\Pgen$, first derived and inferred from sequence
repertoires by Murugan~et~al.~\cite{Murugan2012}, is sharply non-uniform; convergent recombination
makes some sequences enormously over-represented; clonal
expansion and public clones create exact and near duplicates. A query in a common, high-$\Pgen$
region of sequence space has many neighbours \emph{for purely generative reasons}. An i.i.d.\ null
would flag these as wildly significant, which is biologically meaningless. The signal we actually
want is the opposite: \emph{more} neighbours than the background generative process predicts, the
hallmark of antigen-driven selection.

Our approach: define the null by an empirical background distribution $\Pzero$ that carries the
generative and baseline-sharing redundancy but no antigen-driven enrichment, estimate the
null neighbourhood mass from a matched control repertoire (the user-supplied
\texttt{isalgo/airr\_control} set), and calibrate the E-value against it. This is a rigorous,
finite-sample generalization of Karlin--Altschul to a non-i.i.d., biologically structured null, and
the statistical formalization of TCRNET-style neighbour counting against a real
control~\cite{Ritvo2018,Pogorelyy2019}.

\section{Setup and notation}\label{sec:setup}

Let $\Sigma$ be the amino-acid alphabet and $\mathcal{X} = \bigcup_{L\ge 0}\Sigma^L$ the space of
CDR3 sequences. The search engine defines, for a query $q$ and budget $\theta\ge 0$, a non-negative
score $s_\theta(q,x)$ and a \emph{ball}
\begin{equation}\label{eq:ball}
  \ball{\theta}{q} = \{x \in \mathcal{X} : s(q,x) \le \theta\}, \qquad s(q,q)=0,\; s\ge 0.
\end{equation}
The score need not be a metric: with a substitution matrix it is the squared-distance penalty
$\mathrm{pen}(a,b)=s_{aa}+s_{bb}-2s_{ab}$ summed along the optimal alignment (plus gap costs); in
unit-cost mode it is an edit count. Both define a legitimate ball. We work with two background laws
on $\mathcal{X}$: the \emph{realized-repertoire background} $\Pzero$ (what a healthy, unselected
repertoire instantiates) and the \emph{generation law} $\Pgen$ (the V(D)J model). Let
\begin{equation}\label{eq:pi0}
  \pi_0(q,\theta) = \Pzero\!\left(\ball{\theta}{q}\right) = \sum_{x \in \ball{\theta}{q}} \Pzero(x),
  \qquad
  \pi_{\mathrm{gen}}(q,\theta) = \Pgen\!\left(\ball{\theta}{q}\right).
\end{equation}
A control sample $C=(C_1,\dots,C_M)$ and a target set $D$ are given. Crucially, all counts are over
\emph{distinct clonotypes}: the engine deduplicates hits by reference id, so
\begin{equation}
  n_S(q,\theta) = \#\{\,x \in S \text{ distinct} : x \in \ball{\theta}{q}\,\}, \qquad S\in\{C,D\}.
\end{equation}
Write $N=|D|$ (unique clonotypes; see \S\ref{sec:overdispersion}).

\begin{lemma}[Scope monotonicity]\label{lem:monotone}
The balls nest, $\ball{\theta}{q}\subseteq\ball{\theta'}{q}$ for $\theta\le\theta'$ (since $s\ge0$ and
the cut is by a single threshold). Hence $\pi_0(q,\cdot)$, $n_S(q,\cdot)$, the intensity
$\lambda(q,\cdot)$ and the E-value $E(q,\cdot)$ are all non-decreasing in the scope/budget $\theta$,
and the closest-hit score $S_{\min}(q)$ of \S\ref{sec:gumbel} is the smallest $\theta$ with
$n_D(q,\theta)>0$. This justifies sweeping $\theta$ to trace an E-value curve per query.
\end{lemma}

\begin{assumption}[Exchangeability under $H_0$]\label{ass:exch}
Under the null, the unique clonotypes of $D$ are exchangeable with marginal law $\Pzero$.
\end{assumption}
\begin{assumption}[Independent control draws]\label{ass:indep}
The unique clonotypes of $C$ are i.i.d.\ (or exchangeable) $\sim \Pzero$.
\end{assumption}
\begin{assumption}[Background match]\label{ass:match}
$C$ and $D$ share the background $\Pzero$ (same generation $+$ sampling process, matched chain,
species and length composition). All validity is conditional on Assumption~\ref{ass:match}.
\end{assumption}

\section{Null hypothesis and estimator hierarchy}\label{sec:null}

\begin{definition}[Per-query null]
$H_0(q)$: the neighbours of $q$ in $D$ arise from $\Pzero$ with no antigen-driven excess, i.e.\
each $x\in D$ satisfies $\E[\one(x\in\ball{\theta}{q})] = \pi_0(q,\theta)$. The alternative
$H_1(q)$ posits excess mass $\pi_D(q,\theta) > \pi_0(q,\theta)$.
\end{definition}

\begin{lemma}[Self-match exclusion / punctured null]\label{lem:selfmatch}
When the query is itself a database member ($q\in D$, as in a VDJdb-vs-VDJdb scan), the count
$n_D(q,\theta)$ contains the exact self-match (and any exact duplicates of $q$), which are
\emph{deterministic} identity hits, not random draws from $\Pzero$. Including them biases both the
observed count and the null. The correct neighbour statistic is the \emph{punctured} count over the
distance-positive ball,
\begin{equation}\label{eq:punctured}
  n_D^{>}(q,\theta) = \#\{x\in D : 0 < s(q,x) \le \theta\},
\end{equation}
with null intensity $\lambda^{>}(q,\theta) = (N-m_q)\,\pi_0^{>}(q,\theta)$, where $m_q$ is the
multiplicity of $q$ in $D$ and $\pi_0^{>} = \Pzero(\ball{\theta}{q}) - \Pzero(\{x:s(q,x)=0\})$ removes
the point mass at exact matches. The control estimator is punctured identically,
$\hat\pi^{>} = n_C^{>}(q,\theta)/M$, so the deterministic identity term cancels in the calibrated
E-value. (For $q\notin D$ the puncture is vacuous and $n_D^{>}=n_D$.)
\end{lemma}

The estimand is the per-query Poisson intensity $\lambda(q,\theta) = N\,\pi_0(q,\theta)$ (read as
$\lambda^{>}$ with the puncture of Lemma~\ref{lem:selfmatch} whenever $q\in D$). Two
estimators of $\pi_0$ target \emph{different} nulls and must not be conflated.
\begin{itemize}
  \item \textbf{Control / Monte-Carlo (primary):} $\hat\pi(q,\theta) = n_C(q,\theta)/M$, unbiased
  for $\Pzero(\ball{\theta}{q})$, with $M\hat\pi \sim \Binomial(M,\pi_0)$ under
  Assumption~\ref{ass:indep}. It captures the \emph{realized} background, including public-clone
  sharing and finite-repertoire convergence.
  \item \textbf{Generation / analytic (cross-check):} $\hat\pi_{\mathrm{gen}}(q,\theta) =
  \sum_{x\in\ball{\theta}{q}} \Pgen(x)$, computed by enumerating the (small, for small $\theta$)
  ball with the engine and weighting by the V(D)J generation probability of the Murugan~et~al.\
  model~\cite{Murugan2012}. It targets the pure generation null $\Pgen(\ball{\theta}{q})$, which
  omits selection and sampling.
\end{itemize}

\begin{remark}[Selection factor and the thymic correction]\label{rem:thymic}
$\Pgen$ overcounts relative to the naive-repertoire frequency because only a fraction of generated
receptors survive thymic and peripheral selection. Elhanati~et~al.~\cite{Elhanati2014} model this
with a per-sequence \emph{selection factor} $Q(x)\ge 0$, inferred by maximum likelihood and
normalized so $\langle Q\rangle = \sum_x \Pgen(x) Q(x) = 1$, giving the post-selection law
$\Pzero(x) = Q(x)\,\Pgen(x)$. Restricting to a defined cell subset (e.g.\ naive) multiplies $\Pgen$
by a further \emph{global} acceptance fraction; averaged over sequences this thymic factor is
approximately $q \approx 1/2.7$, so an effective pre-immune frequency is
$f(x) \approx q\,Q(x)\,\Pgen(x)$, and baseline antigen-specific T-cell frequencies estimated this way
agree with observed repertoires~\cite{Pogorelyy2018}. Two consequences. (i) The empirical control
$\Pzero$ already \emph{is} the post-selection law, so $\hat\pi$ needs no $Q$; the structured factor
$Q$ matters only for the analytic estimator, where one uses $Q\,\Pgen$ in place of $\Pgen$.
(ii) The global $q$ is sequence-independent and cancels in any ratio or in $\hat\pi$; it matters only
when an \emph{absolute} generation-based frequency is needed (the $\hat\pi_{\mathrm{gen}}$ fallback
for rare queries, \S\ref{sec:precision}).
\end{remark}

\begin{lemma}[The two nulls differ]\label{lem:hierarchy}
In general $\Pzero \ne \Pgen$: thymic and peripheral selection deplete some motifs while
finite-sample public-clone sharing enriches others, so neither $\pi_0 \le \pi_{\mathrm{gen}}$ nor
the reverse holds universally. Hence $\hat\pi_{\mathrm{gen}}$ is used as a variance-reducing control
variate and as a fallback for queries too rare for the control (\S\ref{sec:precision}), not as a
substitute for $\hat\pi$.
\end{lemma}

\section{Poisson approximation with an explicit error bound}\label{sec:poisson}

Fix $q,\theta$. For the unique clonotypes $x_1,\dots,x_N$ of $D$ set
$X_i=\one(x_i\in\ball{\theta}{q})$, $p_i=\E X_i=\pi_0$, $W=\sum_i X_i$, $\lambda=\sum_i p_i = N\pi_0$,
and let $Z\sim\Poisson(\lambda)$. We use the following standard objects. $\mathcal{L}(W)$ denotes the
\emph{law} (probability distribution) of $W$. The \emph{total-variation distance} between two laws
$\mu,\nu$ on $\mathbb{Z}_{\ge 0}$ is
\begin{equation}\label{eq:dtv}
  \dTV(\mu,\nu) = \sup_{A\subseteq\mathbb{Z}_{\ge0}} |\mu(A)-\nu(A)| = \tfrac12\sum_{k\ge0}|\mu(k)-\nu(k)|,
\end{equation}
so a bound on $\dTV$ bounds the error of \emph{every} event probability simultaneously. A family of
\emph{dependency neighbourhoods} is a choice, for each index $i$, of a set $B_i\ni i$ such that $X_i$
is independent of (or nearly independent of) $\{X_j : j\notin B_i\}$; intuitively $B_i$ collects the
clonotypes whose ball-membership is statistically coupled to $x_i$'s (here, those sharing a motif).
The residual $b_3$ below measures exactly how far that near-independence falls short.

\begin{theorem}[Chen--Stein bound \cite{Arratia1989,Arratia1990}]\label{thm:cs}
For any dependency neighbourhoods $\{B_i \ni i\}$,
\begin{equation}\label{eq:cs}
  \dTV\!\big(\mathcal{L}(W),\Poisson(\lambda)\big) \le b_1+b_2+b_3,
\end{equation}
with $b_1=\sum_i\sum_{j\in B_i}p_ip_j$, $b_2=\sum_i\sum_{j\in B_i, j\ne i}\E[X_iX_j]$, and
$b_3=\sum_i \E\big|\,\E[X_i-p_i\mid \sigma(X_j: j\notin B_i)]\,\big|$.
\end{theorem}

\begin{corollary}[Le Cam regime \cite{LeCam1960}]\label{cor:lecam}
If the collapsed clonotypes are independent under $H_0$, take $B_i=\{i\}$; then $b_2=b_3=0$ and
\begin{equation}\label{eq:lecam}
  \dTV\!\big(\mathcal{L}(W),\Poisson(\lambda)\big) \le \sum_i p_i^2 = N\pi_0^2 = \lambda\,\pi_0 .
\end{equation}
The bound is small precisely in the regime of interest: a rare ball $\pi_0\ll 1$ with moderate
$\lambda$ gives error $\le \lambda\pi_0 \to 0$.
\end{corollary}

\begin{corollary}[Void and tail probabilities]\label{cor:void}
With $w=n_D(q,\theta)$ observed,
\begin{align}
  p_{\mathrm{any}}(q,\theta) &= \Prob(W\ge 1) = 1-e^{-\lambda} + O(N\pi_0^2), \label{eq:pany}\\
  p(q,\theta) &= \Prob(Z \ge w) = 1 - \sum_{k<w} \frac{e^{-\lambda}\lambda^k}{k!},
  \quad |\,\Prob(W\ge w) - p(q,\theta)\,| \le b_1+b_2+b_3. \label{eq:ptail}
\end{align}
Both follow from Theorem~\ref{thm:cs}: the void probability is the event $A=\{0\}$ and the tail is
$A=\{w,w+1,\dots\}$, and by~\eqref{eq:dtv} the error on any single event is at most $\dTV\le
b_1+b_2+b_3$ (so $O(N\pi_0^2)$ in the independent regime, Corollary~\ref{cor:lecam}).
\end{corollary}

\begin{remark}[Where biology enters]
Convergent recombination makes the dependency neighbourhood $B_i$ nontrivial: $j\in B_i$ when $x_i$
and $x_j$ share a high-$\Pgen$ motif. The term $b_2=\sum_i\sum_{j\in B_i}\E[X_iX_j]$ is the excess
\emph{pairwise} ball co-occupancy. Under $H_0$ it is controlled by the pairwise ball mass, estimable
from the control by counting \emph{pairs} of control sequences both in $\ball{\theta}{q}$; the
Poisson regime holds when this estimate is $\ll\lambda$. The very same $b_2$ is inflated under $H_1$
(antigen-driven clusters co-occupy the ball), so it is simultaneously the null error term and the
quantity carrying the signal.
\end{remark}

\section{Clonal redundancy and over-dispersion}\label{sec:overdispersion}

\begin{proposition}[Collapsing restores Poisson]\label{prop:collapse}
Let the raw target carry clonotypes with multiplicities (clone sizes) $m_x$. Counting reads/cells in
the ball gives a compound-Poisson total $T=\sum_{k=1}^{K} m_k$ with $K\sim\Poisson(\lambda)$ and
$m_k$ i.i.d.\ $\sim G$, so $\E T=\lambda\mu_G$ and $\Var T=\lambda\,\E[m^2]$, with over-dispersion
index $\Var T/\E T=\E[m^2]/\mu_G\ge 1$. Collapsing to unique clonotypes is the projection $G\equiv1$,
which removes multiplicity-driven over-dispersion and returns the Poisson count $W$ of
\S\ref{sec:poisson}. We therefore deduplicate $C$ and $D$ to unique clonotypes by default.
\end{proposition}

\begin{proposition}[Negative-binomial robustness check]\label{prop:nb}
If multiplicities must be modelled (e.g.\ read-level tests) and $G$ is geometric, $T$ is
negative-binomial; report the NB tail $\Prob(\mathrm{NB}\ge w)$ with mean $\lambda\mu_G$ and
dispersion estimated from observed clone sizes. Under $H_1$ antigen-driven clones are stochastically
larger, so power is retained; the collapsed Poisson test remains the assumption-light default.
\end{proposition}

\begin{proposition}[tf--idf is self-information weighting]\label{prop:tfidf}
Weighting each target hit $x$ by its background self-information $w(x) = -\log\Pzero(\{x\}$-ball$)$
makes the expected per-hit contribution constant under $H_0$; the inverse-document-frequency weight
is exactly the inverse background ball mass, and the term frequency is the clone multiplicity. In the
rare regime the control-set E-value satisfies $E \approx e^{-\sum_{x} \mathrm{idf}(x)}$, so the
``control-set'' and ``tf--idf'' approaches to redundancy are one object.
\end{proposition}

\section{The E-value and multiple testing}\label{sec:evalue}

\begin{definition}[E-value]\label{def:evalue}
For a query family $\mathcal{Q}$, the expected number of background hits is
\begin{equation}\label{eq:evalue}
  E_{\mathrm{tot}}(\theta) = \E_{H_0}[\#\text{hits}] = \sum_{q\in\mathcal{Q}} N\,\pi_0(q,\theta)
  = \sum_{q\in\mathcal{Q}} \lambda(q,\theta).
\end{equation}
The per-query specialization $\mathcal{Q}=\{q\}$ gives the BLAST-convention E-value
$E(q,\theta)=N\pi_0(q,\theta)$, estimated by
\begin{equation}\label{eq:Ehat}
  \widehat E(q,\theta) = \frac{N}{M}\, n_C(q,\theta), \qquad p_{\mathrm{any}} = 1-e^{-\widehat E}.
\end{equation}
\end{definition}

\begin{proposition}[Assumption-free expectation]\label{prop:linearity}
Equation~\eqref{eq:evalue} holds by linearity of expectation regardless of any dependence among
hits (clonal, convergent, or across $\mathcal{Q}$). Consequently $\Prob(\#\text{false hits}\ge 1)\le
E_{\mathrm{tot}}$ by Markov's inequality, and $E_{\mathrm{tot}}$ bounds the expected number of false
discoveries. This robustness---no independence needed for the \emph{mean}---is why the E-value, not
the Poisson tail, is the primary report.
\end{proposition}

\begin{proposition}[Family-wise and false-discovery control]\label{prop:multiple}
Two thresholding regimes for a family of $|\mathcal{Q}|$ tested queries:
\begin{enumerate}
  \item \emph{E-value / Bonferroni (FWER).} Reporting every query with
  $\widehat E(q,\theta)\le\alpha/|\mathcal{Q}|$ controls the family-wise error rate at level $\alpha$:
  by Proposition~\ref{prop:linearity} the expected number of false positives is
  $\sum_q \widehat E \le \alpha$, and $\Prob(\ge 1\text{ false positive})\le\alpha$ by Markov. No
  independence is required. A fixed E-value cutoff (e.g.\ $\widehat E\le 1$, the BLAST default) is the
  $\alpha=|\mathcal{Q}|$ case and bounds the \emph{expected count} of false positives by $1$.
  \item \emph{p-value / Benjamini--Hochberg (FDR).} Using the per-query enrichment p-values
  $p(q,\theta)=\Prob(Z\ge n_D^{>}(q,\theta))$ from Corollary~\ref{cor:void}, the
  Benjamini--Hochberg procedure~\cite{Benjamini1995}---sort $p_{(1)}\le\dots\le p_{(|\mathcal{Q}|)}$,
  reject the $k$ largest with $p_{(k)}\le \tfrac{k}{|\mathcal{Q}|}\alpha$---controls the false
  discovery rate at $\alpha$ under positive dependence of the test statistics, the relevant regime
  here (convergent clusters induce positive correlation).
\end{enumerate}
\end{proposition}

\begin{proposition}[Detectability / minimum cluster size]\label{prop:power}
Under $H_1(q)$ let the antigen-driven excess add $k$ neighbours beyond the background mean
$\lambda=\lambda^{>}(q,\theta)$, so $n_D^{>}\approx\lambda+k$. The enrichment test at E-value cutoff
$\widehat E\le\alpha$ rejects when $n_D^{>}\ge w_\alpha$, the smallest $w$ with
$\Prob(\Poisson(\lambda)\ge w)\le\alpha$. For small $\lambda$ (the typical rare-ball regime),
$w_\alpha$ grows only logarithmically, $w_\alpha \approx \dfrac{\log(1/\alpha)}{\log\log(1/\alpha)-\log\lambda}$
by the Poisson right tail, so a cluster of a handful of convergent neighbours is already detectable;
for moderate $\lambda$ the Gaussian approximation gives the familiar
$k \gtrsim z_{1-\alpha}\sqrt{\lambda}$. The control size enters only through the resolution of
$\widehat\lambda$ (\S\ref{sec:precision}): $M$ must be large enough that the sampling noise of
$\widehat E$ is below the excess $k$ being claimed.
\end{proposition}

\section{How large must the control be?}\label{sec:precision}

Since $M\hat\pi\sim\Binomial(M,\pi_0)$, $\Var\hat\pi=\pi_0(1-\pi_0)/M$ and the relative error is
$\mathrm{CV}(\hat\pi)\approx (M\pi_0)^{-1/2}$.
\begin{proposition}[Resolution]\label{prop:resolution}
To resolve a target E-value $E^\ast=N\pi_0$ to relative error $\rho$ requires
\begin{equation}\label{eq:resolution}
  M \gtrsim \frac{1}{\rho^2\pi_0} = \frac{N}{\rho^2 E^\ast}.
\end{equation}
Resolving $E^\ast\sim 1$ to $10\%$ thus needs $M\sim 100\,N$.
\end{proposition}
\begin{proposition}[Empty-ball regime]\label{prop:zero}
If $n_C=0$, the point estimate $\hat\pi=0$ is degenerate; use the rule of three $\pi_0 \lesssim 3/M$
(95\%), or a $\mathrm{Beta}(n_C+a,\,M-n_C+b)$ posterior, which propagates control uncertainty into a
Poisson--Gamma (negative-binomial) posterior-predictive p-value for $n_D$. The implementation reports
the rule-of-three upper bound $\widehat E \le 3N/M$ when $n_C=0$.
\end{proposition}
When $M$ is inadequate for a rare $q$, the analytic $\hat\pi_{\mathrm{gen}}$ is exact per query for
any $M$ and serves as the fallback (Lemma~\ref{lem:hierarchy}).

\section{Composition and length are handled automatically}\label{sec:composition}

\begin{proposition}\label{prop:composition}
Because $\pi_0(q,\theta)=\Pzero(\ball{\theta}{q})$ is computed for the \emph{specific} query $q$, the
control estimator $n_C(q,\theta)/M$ conditions on $q$'s length and composition automatically: the
same biases that make $q$ common make $n_C$ large. This is the finite-sample, composition-exact
analogue of the Karlin--Altschul $K\,mn$ length normalization, which is needed precisely because the
i.i.d.\ background is query-independent. The only caveat is statistical: rare $q$ require adequate
$M$ (\S\ref{sec:precision}), else fall back to $\hat\pi_{\mathrm{gen}}$.
\end{proposition}

\section{The closest hit: an extreme-value law}\label{sec:gumbel}

\begin{theorem}[Poisson $\Rightarrow$ Gumbel]\label{thm:gumbel}
Let $\lambda(q,t)=N\,\Pzero(\ball{t}{q})$. By the Poisson approximation applied at each radius,
\begin{equation}
  \Prob\!\big(S_{\min}(q) > t\big) \approx e^{-\lambda(q,t)} = e^{-N\Pzero(\ball{t}{q})}.
\end{equation}
If $\log\Pzero(\ball{t}{q}) \approx a + \beta t$ (log-linear ball-mass growth, the generic regime),
then the best score $Y=-S_{\min}$, centred at $u_N=(\log N + a)/\beta$, obeys
$\Prob(Y-u_N\le y) \to \exp(-e^{-\beta y})$, a Gumbel law with scale $1/\beta$. Here $\beta$ is the
\emph{empirical} ball-mass log-slope (regress $\log n_C(q,t)$ on $t$), not the Karlin--Altschul
$\lambda^\ast$. (For lattice scores the Gumbel carries the usual periodic correction.)
\end{theorem}

\section{Relation to Karlin--Altschul}\label{sec:ka}

\begin{theorem}[KA is the product-measure, ungapped case]\label{thm:reduction}
If $\Pzero=\bigotimes_\ell p$ is a product measure and $s$ is the ungapped additive score, then
$\Pzero(\ball{\theta}{q})$ factorizes and, by Cram\'er's theorem, $-\frac1{|q|}\log
\Pzero(\ball{\theta}{q}) \to \lambda^\ast$, the Karlin--Altschul parameter solving $\sum_{ij}p_ip_j
e^{\lambda^\ast s_{ij}}=1$. The intensity $\lambda(q,\theta)=N\Pzero(\ball{\theta}{q})$ then reduces
to $E=K\,m\,n\,e^{-\lambda^\ast S}$ with $K$ the Poisson-clumping constant, recovering
\cite{Karlin1990,Karlin1993,Altschul1997}.
\end{theorem}
Thus the present framework generalizes Karlin--Altschul in three ways: (i) the product measure
$\bigotimes p$ is replaced by the empirical/generative background $\Pzero$; (ii) gaps and
matrix-weighted balls are admitted via the engine's score; (iii) the asymptotic constants
$K,\lambda^\ast$ are replaced by a finite-$N$, finite-$M$ non-asymptotic error bound
(Theorem~\ref{thm:cs}).

\section{The epitope case: a limitation}\label{sec:epitope}

\begin{remark}
For TCR CDR3, $\Pzero$ is generation/repertoire-driven and a healthy-donor control instantiates it.
For \emph{epitopes} (MHC-presented peptides) the relevant background is presentation, not V(D)J
generation. The machinery of \S\S\ref{sec:setup}--\ref{sec:gumbel} applies verbatim with
$\Pzero:=\Pzero^{\mathrm{pep}}$ and a presented-peptide control, but: (L1) there is no closed-form
V(D)J-style generation model $\Pgen$, so only the empirical estimator survives; (L2) presentation is HLA-restricted,
so $\Pzero^{\mathrm{pep}}$ is allele-conditional and the control must be HLA-matched or marginalized
over an HLA frequency distribution; (L3) anchor-residue structure argues for position-weighted ball
geometry. We therefore claim soundness for epitopes only to the extent that a faithful
presented-peptide control is available; no generation-based null is claimed there.
\end{remark}

\section{Practical defaults and algorithm}\label{sec:practice}

\begin{enumerate}
  \item Deduplicate $C$ and $D$ to unique clonotypes (Proposition~\ref{prop:collapse}).
  \item Build a \texttt{seqtree} index of $C$; for each query compute $n_C$ and $n_D$ at scope
  $\theta$ via batched search. When the query may itself be in $D$ or $C$, use the punctured counts
  $n^{>}$ that drop distance-zero (exact/self) hits (Lemma~\ref{lem:selfmatch}).
  \item Report $\widehat E=(N/M)\,n_C^{>}$ (Eq.~\eqref{eq:Ehat}), $p_{\mathrm{any}}=1-e^{-\widehat E}$,
  and $p_{\mathrm{enrich}}=\Prob(\Poisson(\widehat E)\ge n_D^{>})$; use the rule of three when
  $n_C^{>}=0$ (Proposition~\ref{prop:zero}).
  \item Across a query family, threshold on $\widehat E$ for FWER control or apply
  Benjamini--Hochberg to the $p_{\mathrm{enrich}}$ for FDR control (Proposition~\ref{prop:multiple}).
  \item Validate the Poisson regime via the pairwise co-occupancy estimate of $b_2$; if inflated,
  use the negative-binomial check (Proposition~\ref{prop:nb}).
  \item Size the control by Eq.~\eqref{eq:resolution}; fall back to the model-based
  $\hat\pi_{\mathrm{gen}}=\sum_{\ball{\theta}{q}} q\,\Pgen$ (Murugan model, thymic factor
  $q\approx1/2.7$) for rare queries.
\end{enumerate}

This is implemented in \texttt{seqtree.evalues} (with \texttt{exclude\_exact} for the punctured
counts), a thin layer over batched search; the control loader \texttt{seqtree.load\_control} supplies
a deduplicated background.

\bibliographystyle{plain}
\bibliography{refs}

\end{document}
